% !TeX program = xelatex
\documentclass[margin = 5mm, tikz]{standalone}
\usepackage{xfrac}
\usepackage[MnSymbol]{mathspec}
\usepackage{esvect}
\setmainfont[Mapping=tex-text,Ligatures={NoRequired,NoCommon,NoContextual}]{Comic Sans MS}
\setmathfont(Digits,Latin,Greek)[Numbers={Lining,Proportional}]{Comic Sans MS}
\usetikzlibrary{shapes.geometric,decorations,decorations.pathmorphing}
\usetikzlibrary{intersections}
\usetikzlibrary{patterns}
\usetikzlibrary{calc}
\usetikzlibrary{shapes.misc}
\usetikzlibrary{spy}

\pgfdeclarelayer{bg}
\pgfdeclarelayer{bbg}
\pgfsetlayers{bbg, bg, main}

\definecolor{c1}{HTML}{ac0000}
\definecolor{c2}{HTML}{00a0d5}

\newcommand{\abs}[1]{{\left|{#1}\right|}}

\tikzset{
    point/.style n args = {1}{circle, draw = black, inner sep = #1pt, fill = black},
    empty point/.style = {circle, draw = black, inner sep = 2pt, fill = white},
    cross/.style = {cross out, draw, minimum size = 2 * (#1 - \pgflinewidth), inner sep = 0pt, outer sep = 0pt},
    xkcd/.style n args = {1}{decorate, decoration = {random steps, pre length = #1mm, post length = #1mm, segment length = 0.6mm, amplitude = 0.2pt}},
    xkcd/.default = {4},
    point/.default = {2}
}

\begin{document}
    \begin{tikzpicture}
    	\draw[xkcd, thick, black, ->] (-1,0) -- (9,0) node[pos=1,below]{$t$};
    	\draw[xkcd, thick, black, ->] (0,-1) -- (0,9) node[pos=1,left]{$c$};
    	\draw[xkcd, thick, dashed] (8,0) -- (8,8) node[pos=0, below] {$1$} -- (0,8) node[pos=1, left] {$1$};
    	\node (o) at (-0.3,-0.3) {$0$};
    	\draw[xkcd, thick, c1, domain = {0}:{1}, scale = 8] plot (\x, {pow(1 - \x, 3)}) node[right, yshift=-2.5ex]{$c_0(t)=(1-t)^3$};
    	\draw[xkcd, thick, c2, domain = {0}:{1}, scale = 8] plot (\x, {3 * \x * pow(1 - \x, 2)}) node[right, yshift=2ex]{$c_1(t)=3t(1-t)^2$};
    	\draw[xkcd, thick, c1, dashed, domain = {0}:{1}, scale = 8] plot (\x, {3 * pow(\x, 2) * (1 - \x)}) node[right, yshift=5ex]{$c_2(t)=3t^2(1-t)$};
    	\draw[xkcd, thick, c2, dashed, domain = {0}:{1}, scale = 8] plot (\x, {pow(\x, 3)}) node[right]{$c_3(t)=t^3$};
    \end{tikzpicture}
\end{document}